\scanpage{53}
\mspage{2}

est un homotopisme s'il existe \deleted{\unsure{$g$}} 
\[
g:Y\longrightarrow X
\]

telle que
\[
gf\sim\mathrm{id}_X,\quad fg\sim\mathrm{id}_Y \tag*{(\unsure{*|H})}
\]
où la relation $\sim$ désigne l'homotopie dans $\underline{\mathrm{Hom}}(X,X)$ et $\underline{\mathrm{Hom}}(Y,Y)$ respectivement (définie par chemins allant d'un élément à l'autre). Prouvons le

\uline{Lemme} Supposons que $W$ satisfait \textcircled{1} et le a) de la prop. i.e $X\times\Delta^1\longrightarrow X$ dans $W$ pour tt $X$ dans $\mathrm{Cat}$. Alors si \deleted{\unsure{$f,g$}} $X,Y$ dans $\mathrm{Cat}$ et $f,g$ dans $\underline{\mathrm{Hom}}(X,Y)$ sont homotopes, on a
\[
\mathrm{hot}_W(f)=\mathrm{hot}_W(g).
\]
De plus, \deleted{\unsure{on}} $f\in W$ ssi $g\in W$.

Il suffit de le prouver quand il existe un \unsure{hom|lien} $f\longrightarrow g$. Cela peut s'interpréter par l'existence d'une flèche
\[
X\times\Delta^1\xrightarrow{F}Y
\]

dont les restrictions à $X\times\{0\}$ et $X\times\{1\}$ respectivement \deleted{sont}\inserted{soient} $f,g$. \unsure{Pour} prouver que $\mathrm{hot}_W(f)=\mathrm{hot}_W(g)$, il suffit de prouver que
\[
i_0,i_1:X\unsure{\text{---}|\longrightarrow}X\times\Delta^1
\qquad x\longmapsto(x,0),\ x\longmapsto(x,1)
\]
sont égales dans $\mathrm{Hot}_W$, puisque
\[
f=F\circ i_0,\quad g=F\circ i_1.
\]

Or cela résulte du fait que ce sont des sections de $X\times\Delta^1$ sur $X$, et que dans $\mathrm{Hot}_W$ la projection $X\times\Delta^1\xrightarrow{p}X$ est iso. (donc $i_0=i_1=p^{-1}$ dans $\mathrm{Hot}_W$).

\inserted{De plus, le diagramme}
\diagram{\begin{tikzcd}[ampersand replacement=\&,row sep=tiny,column sep=small]
X \arrow[r,"i_0"] \arrow[dr,"f"'] \& X\times\Delta^1 \arrow[d,"F"]\\
\& Y
\end{tikzcd}}
\inserted{où $i_0\in W$, montre que $f\in W$ ssi $F\in W$.}

% slanted
\margin{et de \unsure{même} $g\in W$ ssi $F\in W$, d'où $f\in W\Longleftrightarrow g\in W$.}

Revenons à (*), il vient que dans $\mathrm{Hot}_W$ $gf=\mathrm{id}_X$, $fg=\mathrm{id}_Y$, donc $f$ et $g$ sont des $\mathrm{Hot}_W$-isom. inverses l'un de l'autre. \deleted{Pour \unsure{aller} plus \unsure{loin}} D'autre part, $gf\in W$, $fg\in W$, \deleted{conclure $f\in W$ ceci \unsure{vient} \unsure{du fait} que $W$ est \unsure{fortement} saturé !} puisque $\mathrm{id}_X$, $\mathrm{id}_Y\in W$. Donc on \uline{gagne}, si on a formulé \textcircled{1} de façon assez forte. On a prouvé donc $a\Rightarrow c$.

c) $\Rightarrow$ b) est trivial, puisque si $X$ a un objet final, $X\longrightarrow e$ est un homotopisme.

% archive sheet number: 53
